% 可直接放在主文档的 \end{document} 之前，主文档需加载 booktabs 和 amsmath。
\section{EP/PE 因子结果统计}

样本期为 2020 年 1 月 2 日至 2025 年 12 月 31 日，共 1,455 个交易日。
在行情表原始代码上排除 50 只同时满足 900 开头和 \texttt{.BJ} 结尾的 B 股后，
样本包含 5,966 只股票和 7,089,124 条股票日记录。ST/PT 股票不作筛选，
次新股标记仅用于记录，只有停牌日被主动置为缺失值。

原始因子定义为
\begin{equation}
  \mathrm{EP}_{i,t}
  =\frac{\text{TTM 归母净利润}_{i,t}}{\text{总市值}_{i,t}},
  \qquad
  \mathrm{PE}_{i,t}=\frac{1}{\mathrm{EP}_{i,t}}.
\end{equation}
行情日 $t$ 仅使用截至当日收盘后已经可获得的财务信息，并将因子用于下一交易日。
原始 EP 和 PE 各有 6,719,132 条有效观测，有效率均为 94.78\%。

\begin{table}[htbp]
  \centering
  \caption{原始 EP/PE 描述性统计}
  \label{tab:ep-pe-descriptive}
  \begin{tabular}{lrr}
    \toprule
    统计量 & EP & PE \\
    \midrule
    有效观测数 & 6,719,132 & 6,719,132 \\
    均值 & 0.006932 & $-12.9994$ \\
    标准差 & 0.174816 & 12,694.9046 \\
    1\% 分位数 & $-0.416730$ & $-589.5095$ \\
    中位数 & 0.021173 & 26.2371 \\
    99\% 分位数 & 0.174646 & 752.3711 \\
    负值占比 & 21.36\% & 21.36\% \\
    MAD 截断占比 & 11.05\% & 14.16\% \\
    \bottomrule
  \end{tabular}
\end{table}

最新截面为 2025 年 12 月 31 日，共有 5,204 条有效观测；EP 中位数为
0.01313，PE 中位数为 27.31 倍，负 EP/PE 占比为 28.57\%。全样本中，
PE 的最小值和最大值分别约为 $-3.73\times10^{6}$ 倍和
$5.39\times10^{5}$ 倍，说明 EP 接近零时倒数变换会形成明显长尾。
因此，PE 更适合使用中位数、分位数或去极值结果进行解释，EP 则更适合直接进行
横截面比较。经过逐日 MAD 去极值和 Z 标准化后，EP 和 PE 各保留
6,719,131 条有效结果。本节仅报告描述性统计，不包含回归检验。
